% ANNSIM 2027 draft (gate-model simulators vs IQM Sirius). Anonymous for double-blind review.
% Final version: SCS template (scspaperproc, 5-12 pages); this draft uses plain article for portability.
\documentclass[11pt]{article}
\usepackage[T1]{fontenc}
\usepackage[utf8]{inputenc}
\usepackage[margin=2.5cm]{geometry}
\usepackage{amsmath,amssymb}
\usepackage{booktabs}
\usepackage{graphicx}
\usepackage{url}
\usepackage[hidelinks,pdfauthor={},pdftitle={}]{hyperref}
\newcommand{\tofill}[1]{\textbf{[#1]}}

\title{How Faithful Are Quantum Circuit Simulators? Nine Simulators and Three
Noise Scenarios against a 16-Qubit Star Processor}
\author{Anonymous}
\date{}

\begin{document}
\maketitle

\begin{abstract}
\tofill{Draft: the abstract is written after the results.}
\end{abstract}

\section{Introduction}

Simulators carry most of the work on noisy intermediate-scale quantum
devices~\cite{preskill2018}: processor time is scarce and expensive, so
algorithms are designed, tuned and often evaluated in software, and the device
runs only the final experiment. A simulator is useful only to the extent that
it predicts what the device does, yet its fidelity is usually checked on
synthetic benchmarks~\cite{li2023qasmbench} or on its speed~\cite{jamadagni2024bench},
not on the outcome of the task the simulation stands for. A second gap is the
hardware itself: star-shaped processors, in which every qubit couples to a
central resonator through a dedicated MOVE operation, are recent, and the
software ecosystems support them unevenly~\cite{gnatowski2026routing}.

This paper compares nine freely available circuit simulators against a
16-qubit star processor, IQM Sirius, on a workload whose outcome we can
score: fixed-angle QAOA~\cite{farhi2014qaoa} selecting moves in iterated local
search for the permutation flow shop
problem~\cite{wojewodzki2026gcmm,wojewodzki2026aciids,wojewodzki2027icaisc}.
All simulators receive the same calibration data and three noise scenarios of
increasing detail. We separate the effect of the simulation engine from that
of the compiler, compare simulated and measured outcomes both at the level of
single circuits and at the level of the final schedule, and bound every
comparison by the agreement of the hardware with itself. We also report what
each simulator costs and where it stops working. \tofill{One sentence with the
main finding once the results are in.}

\section{Background}

\paragraph{Workload.}
The workload is quantum neighborhood search for the permutation flow shop
problem on the Taillard instances~\cite{taillard1993}. Iterated local
search~\cite{lourenco2003} moves from a permutation to the best one in an
exponential neighborhood. Choosing that move is a quadratic unconstrained
binary optimization (QUBO) problem~\cite{lucas2014ising} over candidate swaps:
the diagonal holds the makespan change of each swap and a penalty forbids
conflicting pairs. A 20-job permutation needs more variables than the device
holds, so the search splits it into overlapping windows and solves one QUBO per
window with fixed-angle QAOA at depth $p$; the lowest-energy sample of 4096
shots selects the move~\cite{wojewodzki2026aciids,wojewodzki2027icaisc}. Linear
neighborhoods (Adjacent, Fibonacci) give windows of 1 to 16 variables, dense
ones (Dynasearch, Motzkin) up to 15.

\paragraph{Star processor.}
IQM Sirius couples its qubits to a computational resonator. A two-qubit
interaction moves the state of one qubit into the resonator, applies a CZ gate
with the other qubit and moves the state back, so one $R_{ZZ}$ interaction of
QAOA compiles to two CZ and about two MOVE gates. Table~\ref{tab:calib} lists
the calibration that defines the noise models; the compiled circuit holds one
qubit more than the window, the resonator.

\begin{table}[t]
\caption{IQM Sirius calibration snapshot of 23 October 2026 (16 operational
qubits of 24; medians over qubits or loci).}
\label{tab:calib}
\centering
\begin{tabular}{lr}
\toprule
Quantity & Value \\
\midrule
PRX (single-qubit) average gate fidelity & 0.9989 \\
CZ fidelity (qubit--resonator) & 0.9833 \\
MOVE fidelity & 0.9903 \\
Readout error (mean of $0\!\to\!1$ and $1\!\to\!0$) & 0.0175 \\
$T_1$ / $T_2$ of the qubits & 30.0 / 10.9 $\mu$s \\
$T_1$ / $T_2$ of the resonator & 7.5 / 3.1 $\mu$s \\
Gate durations PRX / CZ / MOVE & 40 / 104 / 64 ns \\
\bottomrule
\end{tabular}
\end{table}

\paragraph{Noise channels.}
A depolarizing channel with parameter $\lambda$ replaces the state with the
maximally mixed one with weight $\lambda$; a gate of average fidelity $F$ on
$d$-dimensional systems gives $\lambda = d(1 - F)/(d - 1)$. Amplitude and phase
damping model energy relaxation ($T_1$) and dephasing ($T_2$) during a gate, and
a confusion matrix per qubit models readout errors. A simulator can apply the
channels exactly on a density matrix or sample one Kraus branch per run of the
circuit, which turns a noisy simulation into an average over
trajectories~\cite{dalibard1992mcwf}. The second route costs one statevector
simulation per shot but keeps the memory of a pure state.

\section{Simulators}

Table~\ref{tab:sims} summarizes the programs; the paragraphs below describe
each one and what it offers for noisy simulation at 12 to 17 qubits.

\paragraph{Qiskit Aer.}
Aer is the simulation layer of Qiskit~\cite{javadi2024}. It offers dense
statevector and density-matrix methods, matrix product states and stabilizer
methods, and attaches noise to individual instructions through a noise model.
With a statevector, Aer samples one noisy realization of the circuit per shot,
so it reaches the full 17 qubits of a compiled Sirius circuit. The IQM adapter
for Qiskit~\cite{iqmclient} compiles circuits to the native PRX, CZ and MOVE
gates and builds fake backends from an architecture and an error profile,
which we use to obtain a Sirius model from the calibration snapshot.

\paragraph{Cirq with qsim.}
Cirq~\cite{cirq} is Google's framework for writing and compiling circuits;
qsim is its high-performance C++ statevector simulator. Isakov et
al.~\cite{isakov2021qsim} describe how qsim simulates noisy circuits as quantum
trajectories, choosing the outcome of every channel at random during each run,
and how a delayed inner product cuts the cost of low-noise simulation. Each
repetition re-runs the circuit, so the cost grows with the number of shots.
IQM's Cirq adapter, now part of the IQM client, maps circuits to IQM
architectures.

\paragraph{Qulacs.}
Qulacs~\cite{suzuki2021qulacs} is a C++ simulator with a Python interface built
for speed on small and medium circuits. It models noise as probabilistic gates
(single- and two-qubit depolarizing, dephasing and bit-flip noise) and as
general completely positive maps given by Kraus operators, which covers
amplitude damping, and it also holds density matrices. We reach it through
TKET~\cite{sivarajah2020tket}, whose IQM extension provides a third compiler
for the star topology.

\paragraph{MQT DDSIM.}
DDSIM represents states and operations as decision diagrams, which compress
the redundancy of many circuits~\cite{zulehner2019ddsim}. Its noise-aware mode
simulates decoherence and gate errors either stochastically or
deterministically on decision diagrams for density matrices~\cite{grurl2020noise}.
The documented interface applies amplitude damping, depolarization and phase
flips with one global probability, so per-gate calibrated rates can enter only
as an average.

\paragraph{Qibo.}
Qibo~\cite{efthymiou2022qibo} is a framework for simulation with hardware
acceleration on CPUs and GPUs. In its original description, a circuit that
contains noise channels switches the simulation to density matrices, whose memory grows as $4^n$; with noise it
therefore handles the smaller windows of our workload.

\paragraph{CUDA-Q.}
CUDA-Q~\cite{kim2023cudaq} is NVIDIA's platform for hybrid quantum--classical
programs. Its documentation applies noise models on a density-matrix target on
the CPU; on a machine without an NVIDIA GPU this limits noisy runs to small
windows.

\paragraph{PennyLane.}
PennyLane~\cite{bergholm2018pennylane} centers on differentiable hybrid
programs. Its mixed-state device simulates noise channels on density matrices
and is recommended for small systems; noise models written for Qiskit can be
imported, which eases a like-for-like comparison.

\paragraph{Amazon Braket.}
The Braket SDK~\cite{braket} includes local statevector and density-matrix
simulators; the local density-matrix simulator handles noisy circuits of about
twelve qubits, and larger ones run in a paid cloud service that we do not use.

\paragraph{QuTiP.}
QuTiP~\cite{johansson2013qutip,lambert2024qutip5} simulates open quantum
systems through master equations. Its circuit package simulates circuits at
the level of control pulses with noise~\cite{li2022qutipqip}, a description
below the gate level used by the other programs; we use it on a single small
window as a reference for that level of detail.

\paragraph{Stim, and why it is excluded.}
Stim~\cite{gidney2021stim} simulates stabilizer circuits very fast and is the
standard tool for error-correction studies. QAOA circuits contain arbitrary
rotations and are not Clifford circuits, so Stim cannot run the workload.

\paragraph{Related comparisons.}
Benchmarks of simulation packages measure speed on noiseless
statevectors~\cite{jamadagni2024bench} or provide benchmark
circuits~\cite{li2023qasmbench}. Georgopoulos et al.~\cite{georgopoulos2021noise}
built a noise model from device parameters and compared it with hardware runs.
We add three things: many simulators fed with the same calibration, a star
processor, and a comparison on the outcome of an optimization task.

\begin{table}[t]
\caption{Simulators. Noisy method from the documentation; the last two
columns come from our measurements (stage 1).}
\label{tab:sims}
\centering
\small
\begin{tabular}{llllcc}
\toprule
Simulator & Noisy method & IQM compiler & License & max qubits, noisy & s / circuit \\
\midrule
Qiskit Aer & trajectories, density matrix & qiskit-iqm & Apache-2.0 & \tofill{?} & \tofill{?} \\
Cirq + qsim & trajectories & IQM client (Cirq) & Apache-2.0 & \tofill{?} & \tofill{?} \\
Qulacs & probabilistic gates, Kraus & via TKET & MIT & \tofill{?} & \tofill{?} \\
MQT DDSIM & decision diagrams & --- & MIT & \tofill{?} & \tofill{?} \\
Qibo & density matrix & --- & Apache-2.0 & \tofill{?} & \tofill{?} \\
CUDA-Q (CPU) & density matrix & --- & Apache-2.0 & \tofill{?} & \tofill{?} \\
PennyLane & density matrix & --- & Apache-2.0 & \tofill{?} & \tofill{?} \\
Braket (local) & density matrix & --- & Apache-2.0 & \tofill{?} & \tofill{?} \\
QuTiP & master equation (pulses) & --- & BSD-3 & \tofill{?} & \tofill{?} \\
Stim & stabilizer only & --- & Apache-2.0 & n/a & n/a \\
\bottomrule
\end{tabular}
\end{table}

\section{Methodology}

\paragraph{Noise scenarios.}
All simulators use the same calibration snapshot and three scenarios. In S1
the circuit runs without noise. In S2 every gate carries a depolarizing
channel matched to its measured fidelity and every qubit a readout confusion
matrix. S3 adds amplitude and phase damping during each gate. On Sirius,
relaxation alone exceeds the measured error of most gates (median CZ error
0.0167 measured against 0.0205 implied by $T_1$ and $T_2$), so S3 bounds the
noise from above. Neither scenario models idle qubits; every simulator applies
noise during gates only. DDSIM receives S2 and S3 with averaged rates.

\paragraph{Two tracks.}
In track A every simulator runs the same compiled circuit, exported as
OpenQASM with MOVE written as a swap between a qubit and the resonator, so any
difference comes from the engine and its channels. In track B the same logical
circuit passes through three compilers (the IQM adapters for Qiskit and Cirq
and TKET) and runs in one engine, so differences come from compilation.

\paragraph{Hardware references.}
Table~\ref{tab:hw} lists the device data. Run-level comparisons use iterated
local search on Sirius at $p = 1$ from an earlier campaign~\cite{wojewodzki2026aciids}
and at $p = 2, 3$ \tofill{H2}. Circuit-level comparisons need histograms of
single windows; a small run on two different days \tofill{H3} gives them and
measures how well the device agrees with itself.

\paragraph{Metrics.}
For single windows we compare output distributions by total variation distance
and Hellinger distance and report the probability of sampling the optimal move
and whether the lowest-energy of 4096 samples is optimal. Two independent
4096-shot samples of one distribution never agree exactly; we estimate that
floor by bootstrap and the hardware floor from the two days of H3. For whole
runs we compare the relative deviation of the makespan from the Taillard lower
bound after the same number of moves as the device made, with a paired sign
test and Kendall's $\tau$ over instances. Cost is wall time per circuit and
peak memory.

\section{Experimental Setup}

The study runs in three stages. First, every simulator is installed and
tested on one window. Second, all simulators that can run a scenario do so on
the windows of instances 0--3 of the Taillard classes with $n = 20$ jobs: 161
windows with an improving move, of which 71 hold at most 11 variables and fit
the density-matrix simulators. Third, the two or three simulators that best
predict the device at acceptable cost run the full Taillard set ($n = 20$ to
$500$), against the $p = 1$ hardware data available for all twelve classes. We
fix the selection criteria before the second stage.

\begin{table}[t]
\caption{Hardware data used as reference.}
\label{tab:hw}
\centering
\small
\begin{tabular}{llll}
\toprule
Device & Experiment & Scope & Status \\
\midrule
IQM Sirius & ILS, $p = 1$, 10 s, 4096 shots & 12 classes ($n = 20$--$500$), inst.\ 0--3; 225 runs & done (7--14 Oct 2026) \\
IQM Sirius & ILS, $p = 2, 3$ (H2) & $n = 20$, inst.\ 0--3; 96 runs & \tofill{planned} \\
IQM Sirius & window histograms (H3), two days & tai20, inst.\ 0, $p = 1, 2$ & \tofill{optional} \\
IBM Heron (ibm\_fez) & one move per run, $p = 1$ & tai20\_5, inst.\ 0--9 & done (Aug 2026) \\
IQM Garnet, Emerald & one move per run, $p = 1$ & tai20\_5, inst.\ 0 & done (Sept--Oct 2026) \\
\bottomrule
\end{tabular}
\end{table}

\section{Results}

\tofill{Draft: written after stages 1--3.}

\begin{table}[t]
\caption{Track B: gates per window after compilation (median), $p = 1$.}
\label{tab:compilers}
\centering
\small
\begin{tabular}{lccc}
\toprule
Neighborhood & qiskit-iqm (CZ / MOVE) & IQM Cirq (CZ / MOVE) & TKET (CZ / MOVE) \\
\midrule
Adjacent & 72 / \tofill{?} & \tofill{?} & \tofill{?} \\
Fibonacci & 15 / \tofill{?} & \tofill{?} & \tofill{?} \\
Dynasearch & 164 / \tofill{?} & \tofill{?} & \tofill{?} \\
Motzkin & 134 / \tofill{?} & \tofill{?} & \tofill{?} \\
\bottomrule
\end{tabular}
\end{table}

\begin{figure}[t]
\centering
\fbox{\parbox[c][3.5cm][c]{0.9\textwidth}{\centering \tofill{Figure: total
variation distance to Sirius, simulator $\times$ scenario (heat map), with the
shot-noise floor and the hardware day-to-day floor.}}}
\caption{Agreement of each simulator and scenario with the device on single windows.}
\label{fig:heat}
\end{figure}

\begin{figure}[t]
\centering
\fbox{\parbox[c][3.5cm][c]{0.9\textwidth}{\centering \tofill{Figure: relative
deviation of the makespan, simulator versus Sirius, one point per instance, $p = 1, 2, 3$.}}}
\caption{Run-level agreement with the device.}
\label{fig:prd}
\end{figure}

\begin{figure}[t]
\centering
\fbox{\parbox[c][3.5cm][c]{0.9\textwidth}{\centering \tofill{Figure: wall time per
circuit and peak memory versus number of qubits, all simulators (log scale).}}}
\caption{Cost of noisy simulation.}
\label{fig:cost}
\end{figure}

\section{Discussion}
\tofill{Draft: which simulator for which purpose; limitations: resonator modeled
as a qubit, no idle noise, one calibration snapshot, $n = 20$ core.}

\section{Conclusions}
\tofill{Draft.}

\bibliographystyle{plain}
\bibliography{references}
\end{document}
